\documentclass[../../main.tex]{subfiles}
\begin{document}

{\bf [解]}
$n$次多項式を$P_n(x)$とおく．題意（$\diamondsuit$とおく）を帰納的に示す．

まず$n=1$の時は$P_1(x)=ax+b\ (a\ne0)$とおけて，$P_1(0),P_1(1)\in\mathbb{Z}$だから
\begin{align}
a+b\in\mathbb{Z}, \quad b\in\mathbb{Z} \quad\therefore\ a,b\in\mathbb{Z}
\end{align}
となり，たしかに任意の整数 $k$に対し$P_1(k)\in\mathbb{Z}$となって$\diamondsuit$は成立する．

そこで以下$n \le \ell \in\mathbb{N}$での$\diamondsuit$の成立を仮定し，$n=\ell+1$での成立を示す．
$P_{n+1}(x)$について，$P_{n+1}(0),\cdots,P_{n+1}(n+1)$が全て整数だとする．
新しく
\begin{align}
 Q(x)=P_{n+1}(x+1)-P_{n+1}(x) \label{eq:1}
\end{align}
とおくと，$Q(x)$は$n$次以下の多項式であってかつ$Q(0),Q(1),\cdots,Q(n)$が全て整数である．従って帰納法の仮定から，任意の$k\in\mathbb{Z}$に対して
\begin{align}
Q(k)\in\mathbb{Z}
\end{align}
が成り立つ．$P_{n+1}(0),\cdots,P_{n+1}(n+1)$が全て整数だから，\cref{eq:1}を繰り返し用いれば任意の$k\in\mathbb{Z}$に対して$P_{n+1}(k)\in\mathbb{Z}$となり，$n=\ell+1$でも$\diamondsuit$は成立する．

以上より数学的帰納法により題意は示された．

\newpage

\end{document}
